
\examquestion{2025}{3}{Two-photon NOON input and beam-splitter output probabilities}{4}
With the paper's balanced beam splitter
\(c=(a+b)/\sqrt2,\ d=(a-b)/\sqrt2\), the inverse
relations for creation operators are
\[
a^\dagger=\frac{c^\dagger+d^\dagger}{\sqrt2},
\qquad b^\dagger=\frac{c^\dagger-d^\dagger}{\sqrt2}.
\]
\subsection*{(a) Transform the plus NOON input}
Using \(\ket{2,0}_{ab}=(a^\dagger)^2\ket0/\sqrt2\),
\[
\ket{2,0}\ \longrightarrow\
\frac12\ket{2,0}_{cd}
+\frac1{\sqrt2}\ket{1,1}_{cd}
+\frac12\ket{0,2}_{cd}.
\]
Similarly
\[
\ket{0,2}\ \longrightarrow\
\frac12\ket{2,0}_{cd}
-\frac1{\sqrt2}\ket{1,1}_{cd}
+\frac12\ket{0,2}_{cd}.
\]
The initial state is
\(\ket{\Psi_{\rm in}}=(\ket{2,0}+\ket{0,2})/\sqrt2\).
Add the two output \emph{amplitude vectors}:
\[
\boxed{\ket{\Psi_{\rm out}}=
\frac{\ket{2,0}_{cd}+\ket{0,2}_{cd}}{\sqrt2}}.
\]

\subsection*{(b) Detection probabilities}
Read probabilities from squared moduli:
\[
\boxed{P_{20}=\frac12,\qquad
P_{02}=\frac12,\qquad P_{11}=0}.
\]
They add to 1.

\subsection*{(c) Physical interpretation}
The two input components are \emph{coherent}
number-state amplitudes, so the histories leading
to one photon in each output port interfere
destructively. The input is a two-photon path-entangled
NOON state. The output has perfect two-photon bunching,
but here it originates from the plus-NOON superposition:
this is not the same initial state as the ordinary
two-single-photon HOM experiment. The \emph{minus}
NOON state in the same convention would instead
map to \(\ket{1,1}_{cd}\).

\subsection*{(d) An alternative superposition}
For
\(\ket{\Phi_{\rm in}}=(\ket{1,1}+\ket{2,0})/\sqrt2\),
first transform the \(\ket{1,1}\) part:
\[
a^\dagger b^\dagger\ket0
=\frac{(c^\dagger)^2-(d^\dagger)^2}{2}\ket0
=\frac{\ket{2,0}-\ket{0,2}}{\sqrt2}.
\]
This contains \emph{no} \(\ket{1,1}_{cd}\).
The \(\ket{2,0}\) input contributes
\(\ket{1,1}\) with amplitude \(1/\sqrt2\).
Including the initial superposition coefficient
\(1/\sqrt2\) yields
\[
\boxed{\mathcal A_{11}=\frac12,\qquad
P_{11}=|\mathcal A_{11}|^2=\frac14}.
\]
\solcheck{Never square each path probability before
adding coherent amplitudes; the cancellation in
part (a) would be lost. Account for \(\sqrt{2!}\)
when normalising the two-photon states.}
